\section{Conclusion and Future Work}
\label{sec:conclusion}

In this paper, we introduced \textbf{Lyapunov-Stable Active Representation Coupling (L-ARC)}, a training-free dynamic serving framework centered on \emph{state-shielded continuous kinetics}. By modeling representation error as a system-level candidate Lyapunov function, we analytically derived a discrete-time Lyapunov difference equation, yielding a local \textbf{Dissipation Guard} that gates and adapts feedback coupling on-the-fly. Together with our \textbf{Entropy-Gated Concentration Gating (ECG-Reset)} state-shielding filter, L-ARC provides dual-layered representation and state-space shielding.

Our empirical results on the 14-layer Coordinate Sandbox (ICS) demonstrate that L-ARC establishes robust dynamic serving under clean, transient routing failure, and biased routing conditions. We report with complete scientific transparency that while active feedback warping is practically redundant under clean serving ($p=0.106$) and introduces a latency overhead, our proposed \textbf{Entropy-Triggered Gating (ET-L-ARC)} optimization dynamically bypasses Dissipation Guard calculations under low routing uncertainty. This collapses the relative latency overhead to only \textbf{99.85\%} ($0.06$ ms of absolute wall-clock latency per sample), making the system exceptionally practical for edge hardware. Under transient routing failures (Setting C) and confident systematic router bias (Setting D), L-ARC's multi-layered architecture is exceptionally valuable: ECG-Reset provides essential state-space shielding under failures (+5.14\% absolute accuracy gain), our Dissipation Guard provides vital representation-space shielding, and our Representation-Agreement State Correction (RASC) mechanism completely resolves the state-locking failure under persistent router bias (+5.32\% accuracy gain compared to Decay-ChemMerge) to protect semantic representation quality. By comparing L-ARC against simpler smoothing and feedback heuristics, we isolate and demonstrate the irreplaceable role of control-theoretic state-shielding and closed-loop gating under system-level faults.

\subsection{Broader Implications and Future Directions}
The success of L-ARC highlights a broader, powerful design philosophy: **deep learning heuristics can be systematically stabilized and optimized by applying classical control-theoretic principles.** Rather than relying on computationally intensive offline hyperparameter sweeps or training complex parametric gating networks, control-theoretic closed-loop systems can adapt serving parameters dynamically on-the-fly with negligible computational overhead.

Several promising avenues exist for extending this research:
\begin{enumerate}
    \item \textbf{Scaling to Full-Scale Transformers:} While the Analytical Coordinate Sandbox (ICS) serves as an elegant, mathematically controlled testbed for isolating representation flow dynamics, a vital next step is scaling L-ARC to high-dimensional attention states in full-scale transformers (e.g., LLaMA, RoBERTa, or ViT) evaluating standard benchmarks (e.g., GLUE, MMLU). This will allow analyzing the impact of non-linear feedforward blocks on the Layer-Identity Approximation error bounds.
    \item \textbf{Adaptive and Optimal Threshold Tuning for ECG-Reset:} In this work, we introduced the Entropy-Gated Concentration Gating (ECG-Reset) mechanism to shield stateful kinetics from transient routing failures. While our static threshold ($\theta_H = 0.95$) successfully shields the ensembling states, future work could investigate dynamically tuning the entropy threshold based on running statistics or reinforcement learning, optimizing the trade-off between kinetics memory retention and instant error-recovery.
    \item \textbf{Token-Level Lyapunov Routing:} Applying L-ARC to token-level routing in large-scale autoregressive Mixture-of-Experts (MoE) architectures, where token routing noise and representation shift are highly prevalent.
    \item \textbf{Stochastic Lyapunov Control:} Deriving stochastic Lyapunov functions and control bounds to provide probabilistic convergence guarantees in extremely noisy or adversarial serving environments.
    \item \textbf{General PEFT Merging:} Extending Lyapunov-stable feedback loops to multi-modal and cross-attention adapters, enabling stable unified routing across heterogeneous visual, linguistic, and auditory modalities.
    \item \textbf{Non-linear Lyapunov Functions:} Deriving and proving convergence under non-linear Candidate Lyapunov functions that directly account for deep neural network non-linearities, bypassing the Layer-Identity Lipschitz approximation completely.
    \item \textbf{Online Centroid Adaptation under Non-Stationary Environments:} L-ARC relies on early-stage centroids ($\mu_k^{(3)}$) extracted during calibration to compute similarity error. In highly non-stationary environments where downstream domain shifts or coordinate drifts affect representation spaces, these centroids can become biased, challenging RASC's decoupling guarantees. Future work will investigate extending L-ARC with online clustering or EMA-based centroid tracking to adapt centroids on-the-fly, ensuring that our closed-loop decoupling guarantees remain robust under long-term non-stationary shifts.
\end{enumerate}
We hope this work inspires further research into bridging control theory and deep learning serving, leading to mathematically certified, adaptive, and highly robust AI serving systems.
